\section{A fresh e133 host and four complementary original families}
\label{sec:cubic-e133-actual-model}

The missing sextic certificate AC-G131 permits a new construction at
the residual exponent \(133\).  This exponent is not the exponent
\(267\) of AC-G123: their rational residual characters differ by the
quadratic character at \(269\).  We construct a new ordinary branch,
new finite-torus summand, new level-raising prime, and new native
modular ideals.  The comparison with the original models is then
made on their own selected local components.

\paragraph{Original domain and retained inputs.}
Put
\[
 p=269,\quad K=\mathbb Q(\sqrt{-3}),\quad
 \psi_0=\chi_{-4},\quad e_*=133,\quad j\in\{1,2\},
 \quad \bar\gamma_*=\bar\psi_0\bar c_p^{133}.
 \tag{G132.1}
\]
The actual cyclotomic character satisfies
\(c_p(\operatorname{Fr}^{\mathrm{geom}}_w)=q_w^{-1}\).
The cubic character is the fixed character of
\(K(\sqrt[3]{269})/K\).
For the core result take \(N=4\), \(\psi=\psi_0\), and every
positive \(e\equiv133\pmod{268}\).  Keep the original coefficient
DVR \(\mathcal O_0\), fraction field \(E_0\), residue field \(k_0\),
actual normalized cubic ring \(S_j\), selected local prime \(P_*\),
and model
\[
 (A_0,\mathcal A,x,q_0:A_0\longrightarrow S_j,\mathcal Q),
 \qquad I_0=\mathcal Q.\mathrm{ideal}.
 \tag{G132.2}
\]
Write \(\mathscr B_0\) for its complete degree-zero ring and \(L_j\)
for its actual weight-minus-one cotangent length.  The rational
Eisenstein auxiliary prime is denoted \(q_{\rm aux}=3229\), so it
cannot be confused with the ring map \(q_0\).

Retain the literal original \(\mathrm{GD}/\mathrm{Dat}\), coefficient,
marked-element, selected-component and local hypotheses of
AC-G65--AC-G68; \(\mathrm{CohInp}\), \(\mathrm{LemmaP}\), the generic
\(\mathrm{LemmaM},\mathrm{LemmaM'}\), and the original
\(\mathrm{AutomorphyUnramifiedCoefficientsAlt}\) supply.
In particular retain \(e_{\mathcal O_0}=1\) where that original
modularity input requires it.  Retain the published attachment,
ordinary control, transfer, local--global compatibility and cyclic
base-change inputs used in AC-G119--AC-G121; the native
\(\mathrm{LevEx}\), \(\mathrm{ProModularRingExists}\),
\(\mathrm{LamEx}\), and \(\mathrm{ModU}\) supplies; and the adjoint
theorem and exact original-field descent/counting inputs of
AC-G110--AC-G112.  Class field theory, arithmetic duality and the
abelian-field Iwasawa inputs retain their audited scopes in
AC-G74, AC-G76 and AC-G90--AC-G100.
No original Type II endpoint, additional support ideal, mixed
Type II Selmer vanishing, or \(\mathrm{RigiditySecondCase}\) is assumed.

\begin{theorem}[AC-G132-e133-model: fresh actual support and complementary families]
\label{AC-G132-e133-model}
Under these retained inputs there are a finite CM host \(H/F\),
\(F=H^+\), full local groups, a finite host coefficient DVR
\(\mathcal O\), and actual normalized tuple rings
\[
 B\subset A\subset R,\qquad
 \theta=(\bar\chi^j\bar\gamma_*,1,
                         \bar\chi^j\bar\gamma_*^{-1}).
 \tag{G132.3}
\]
They may be chosen once for fixed \(\mathcal O_0,j\) and all
positive \(e\equiv133\pmod{268}\).  They are newly constructed data,
not the rings or modular ideal of AC-G123.
The newly constructed native ideal \(J_3\subset\mathcal B_3\)
and map \(\Phi_3:\mathcal B_3\to B\) give ideals
\[
 I=\Phi_3(J_3)B,\qquad I_{\rm off}=I\cap M_E
 \tag{G132.4}
\]
such that
\[
 \begin{aligned}
 I&\subset\mathfrak q\cap B
       &&\text{for every Type I prime of }A,\\
 I_{\rm off}&\subset\mathfrak q\cap B
       &&\text{for every off-triangular prime of }A,\\
 \sqrt{I_{\rm off}+\lambda_\theta}&=\mathfrak m_B .
 \end{aligned}
 \tag{G132.5}
\]
Here \(\lambda_\theta\) is the actual joint-residual-inertia-kernel
weight ideal, including the host uniformizer.
Every original model (G132.2) has \(\mathscr B_0\) finite over
\(\mathcal O_0\), satisfies \(\mathrm{M1}\)--\(\mathrm{M5}\), and
has only rigid Type I two-sided characteristic-zero local points.

The same conclusions hold on this newly constructed host for all
four rows
\[
\begin{array}{c|c|r|c|c}
(s,t)&\psi_{s,t}=\psi_0\chi_5^s\chi_{269}^t
 &N_{s,t}&e\pmod{268}&\operatorname{comp}_p\psi_{s,t}\\ \hline
(0,0)&\chi_{-4}&4&133&1\\
(1,0)&\chi_{-20}&20&133&1\\
(0,1)&\chi_{-1076}&1076&267&\chi_{269}\\
(1,1)&\chi_{-5380}&5380&267&\chi_{269}
\end{array}
\tag{G132.6}
\]
with their own original selected local primes and original rings.
Thus \(e\equiv133-134t\pmod{268}\); the two rows ramified at \(p\)
are checked with their actual quadratic phase, not by identifying
their characteristic-zero local deformation conditions.

For any one row put
\[
\begin{split}
 \alpha_j&=\chi^{-j}\psi_{s,t}^{-1}c_p^{-e}|_{G_K},\qquad
 \mathcal D_j=(E_0/\mathcal O_0)(\alpha_j),\\
 \mathcal S_j&=\ker\bigl(H^1(G_K,\mathcal D_j)\\
 &\hspace{3em}\longrightarrow
 \prod_{v\mid p}H^1(G_{K_v},\mathcal D_j)
 \times\prod_{v\nmid p}H^1(I_v,\mathcal D_j)\bigr),\\
 s_j&=\operatorname{length}_{\mathcal O_0}(\mathcal S_j)/
             e_{\mathcal O_0},\qquad
 Y=N_{s,t}(e+1),\qquad b=121^e+1+121^{-e}.
\end{split}
\tag{G132.7}
\]
These lengths are finite.  For geometric Frobenius \(\sigma\)
above the inert rational prime \(11\), there is a finite family
\(X_j\) of actual Type I representations, pairwise nonisomorphic
on \(G_K\), with
\[
 s_j\le L_j/e_{\mathcal O_0}
 \le2\sum_{\rho\in X_j}v_{269}(\operatorname{tr}\rho(\sigma)-b),
 \qquad \#X_j\le C_1Y^{A_1}.
 \tag{G132.8}
\]
The constants are precisely those of AC-G112-count.
For a nonempty \(J_0\subset\{1,2\}\), one nonzero rational integer
\(\mathcal N_{s,t,e,J_0}\) satisfies
\begin{align}
 \sum_{j\in J_0}L_j/e_{\mathcal O_0}
       &\le2v_{269}(\mathcal N_{s,t,e,J_0}),\notag\\
 \log|\mathcal N_{s,t,e,J_0}|
       &\le5C_1\#J_0\,Y^{A_1+1}\log11,\notag\\
 \left(\sum_{j\in J_0}L_j/e_{\mathcal O_0}\right)\log269
       &\le10C_1\#J_0\,Y^{A_1+1}\log11.
 \tag{G132.9}
\end{align}
The same bounds hold with \(s_j\) in place of the normalized
cotangent lengths.  The two labels may use separate hosts; the
integer uses their common original-field classifier.
\end{theorem}

\begin{proof}
\proofstep{1. The new finite four-sector certificate.}
Set
\[
 F_1=K\mathbb Q(\sqrt{-23}),\quad E=\mathbb Q(\sqrt5),\quad
 H_{\rm ab}=F_1E,\quad d_{\rm ab}=[H_{\rm ab}^+:\mathbb Q]=4.
 \tag{G132.10}
\]
The four characters of \(H_{\rm ab}/K\) are the restrictions of
\(1,\chi_{69},\chi_5,\chi_{345}\).  For
\(\alpha=\bar\chi^{-j}\bar\psi_0^{-1}\bar c_p^{-133}\),
the exact full-local Shapiro decomposition of AC-G103 gives
\[
 \mathcal S(H_{\rm ab},\alpha)
   =\bigoplus_{\delta=1,\chi_{69},\chi_5,\chi_{345}}
                         \mathcal S(K,\alpha\delta).
 \tag{G132.11}
\]
This includes the actual away conditions at ramified places.
Since \(\bar c_p^{134}=\bar\chi_{269}\), the Kummer character
is
\(\eta_\delta=(\alpha\delta)/\bar c_p
 =\bar\chi^{-j}\bar\psi_0\bar\chi_{269}\delta\).
Its inverse, not itself, is the unit/class eigenspace.
Writing \(t_0^3=269\), its splitting field is
\(M_D=K(t_0,\sqrt D)\), and its real reflection field is
\(L_D=\mathbb Q(u)\), \(u^6=m^3\cdot269\), \(D=269m\):
\[
\begin{array}{c|r|r|c|c}
\delta&m&D&\text{class input}&\dim_k\mathcal S(K,\alpha\delta)\\\hline
1&3&807&\text{AC-G115}&0\\
\chi_{69}&23&6187&\text{certified quotient of }C_{84}\times C_2&\le1\\
\chi_5&15&4035&\text{AC-G126}&0\\
\chi_{345}&115&30935&\text{AC-G131}&\le1
\end{array}
\tag{G132.12}
\]
For \(D=6187\), the preserved successful PARI input uses
\(u^6-3272923\).  Its output is
\(\mathrm{CLASS\_QUOTIENT\_CERTIFY}=1\), followed by
\(\mathrm{DONE}\) and exit code zero.  Under the saved
\texttt{bnfcertify} flag-one specification, the actual class group
is a quotient of \(C_{84}\times C_2\); this does not assert class
group equality or a complete unit basis.  Its \(269\)-part is zero.
AC-G131 proves the same \(269\)-vanishing for \(L_{30935}\)
by exact unit saturation and an isolated rigorous analytic interval,
without a successful PARI class-group certification.

The \(S_3\times C_2\) reflection argument of AC-G103 transfers
these two real class assertions to the nontrivial cubic sector of
\(\operatorname{Cl}(M_D)[269]\).  Equivalently, the order-two
reflection has no invariants in that sector and hence acts as
minus one; its conjugates do so too, forcing their order-three
products to act trivially.  The target nontrivial cubic sector
therefore vanishes.  Localizing at \(p\) cannot introduce a class
sector.  In these nontrivial cubic/quadratic unit sectors the
global unit multiplicity is one, and the \(p\)-prime permutation
correction has no cubic component.  Thus AC-G103's exact formula
\(1+c_\delta-\ell_\delta\) gives the two upper bounds by one.
The zero bounds are the complete strict local calculations of
AC-G115 and AC-G126, for both \(j\); their coefficient-prime
conditions are zero in full local cohomology, not only valuation
conditions.  Consequently
\[
 a_{{\rm ab},0}\le2<4=d_{\rm ab},\qquad r_{\alpha,{\rm ab}}\le2 .
 \tag{G132.13}
\]
This retains every quadratic host factor.  It is not
prime-to-\(p\) rank ascent from the field \(K(t_0)\).

\proofstep{2. A new actual auxiliary seed at weight \(267\).}
Apply the all-odd-weight theorem AC-G119 with \(e_*=133\).
It supplies an actual ordinary cuspidal newform \(f_*\) of weight
\(267\), nebentype \(\eta_K\), primitive level
\(48q_{\rm aux}^2\), and
\[
 V=R_{f_*}c_p^{133},\quad \det V=\eta_K,\quad
 \bar V^{\rm ss}=\eta_K\bar\gamma_*\oplus\bar\gamma_*^{-1}.
 \tag{G132.14}
\]
The source representation \(R_{f_*}\) uses geometric Frobenius;
this is the same normalization as AC-G119.  Its exact extra
inertia at \(q_{\rm aux}\) is
\(\Theta^{1/2}\oplus\Theta^{-1/2}\), with \(N=0\).
The real cyclic degree-\(p\) field \(T_\Theta/\mathbb Q\)
cut out by \(\Theta\) kills this inertia.  It is totally
ramified only at \(q_{\rm aux}\), and \(p\) is unramified and
inert in it; \(\Theta(p)=\zeta_p^{175}\).
The actual ordinary flag at \(p\) has the finite phase
\(\psi_0|_{G_{\mathbb Q_p}}=1\).
No original \(\mathrm{Dat.IsPt}\) is being asserted.

The residual character \(\bar\gamma_*\) has order \(268\) on
coefficient-prime inertia, since \(\gcd(133,268)=1\).
In particular \(\bar\gamma_*^4\ne1\).
If the irreducible \(V\) had a nontrivial quadratic self-twist,
comparison of its two distinct residual constituents would
interchange them and make that twist \(\eta_K\bar\gamma_*^2\);
squaring would give \(\bar\gamma_*^4=1\), a contradiction.
It has no nontrivial finite abelian self-twist of any order:
determinants force the order to divide two.
The character decomposition of the restriction commutant
therefore proves absolute irreducibility on every abelian
initial field in (G132.10), with \(T_\Theta\) also adjoined.
At subsequent cyclic \(p\)-steps it persists because a
self-twist of order \(p\) would have square one.
The same residual argument applies to the classical members
of the new Hida branch.

\proofstep{3. The new branch, exact types, and full unit torus.}
Take an actual ordinary branch through \(f_*\), with finite
torsionfree domain \(R_{\rm H}\) over its fixed Teichmuller
weight disc \(\mathcal O[[W]]\).  Normalize its determinant to
\(\eta_K\).  This branch is not the branch chosen through the
weight-\(535\) point of AC-G120.
Its tame diamond eigenvalues are constant even if their orders
are divisible by \(p\): if \(a^d=1\) and \(a\) specializes to
\(\zeta\), then
\((a-\zeta)Q(a)=0\), while
\(Q(a)(f_*)=d\zeta^{d-1}\ne0\) in characteristic zero.
The domain property forces \(a=\zeta\); no nonintegral
diamond projector is used.

The exact branch local-type theorem used in AC-G120, including
its purity and monodromy clause at the classical seed, applies
to this new branch.  In particular it fixes the displayed
order-\(p\) inertia and \(N=0\) at \(q_{\rm aux}\).
Restrict to
\[
                  F_{\rm pre}=F_1T_\Theta
 \tag{G132.15}
\]
before imposing its hyperspecial level there.  This kills the
actual extra inertia, not just its reduction.
At every other away place the exact rank-two type is
\(\psi_0\oplus\psi_0^{-1}\) after restriction to the CM
field, with \(N=0\).  At \(3\), \(\chi\) is trivial because
\(269\equiv-1\pmod9\) is a cube in \(\mathbb Q_3\);
the rational \(\eta_K\) factor is killed on \(G_K\).
The finite types subsequently averaged have order prime to
\(269\).

The field \(F_{\rm pre}/F_{\rm pre}^+\) is finite-unramified
and relatively split at \(2,3,p,q_{\rm aux}\).
These are the same direct quadratic splitting calculations for
\(\mathbb Q(\sqrt{-23})\) as in AC-G120, independent of the
weight.  Its coefficient-prime completion is the unramified
degree-\(538\) extension of \(\mathbb Q_{269}\).
This is a local number field degree, not a coefficient-ring
ramification index.
The endoscopic starting packet of the retained Thorne theorem
has the required repeated parallel infinity types: over the
quadratic auxiliary CM field their repetition is \(2p\),
which is even.  The actual normalized rank-two input is
cuspidal on \(F_{\rm pre}\) by Step 2.

Use the new symmetric weight
\(\mu_*=(132,0,-132)\).
The full finite unit-torus character is the restriction of
\[
        (\chi^j,1,\chi^j)\,w_0(\mu_*)^{-1}.
 \tag{G132.16}
\]
The full-torus Jacquet calculation of AC-G79 applies here,
not only its pro-\(p\) part.  At each \(p\)-place the ordinary
slopes at a weight \(w\equiv133\pmod{268}\) are
\((-fw,0,fw)\), \(f=538\); the normalized positive operator
has valuation
\(2f-2fw+2f(w-1)=0\), and every other ordering has strictly
positive rearrangement excess.  The parallel Teichmuller
weight difference is a multiple of \(268\), so the dense
classical family lies in the one character (G132.16).
From the start retain the three commuting integral projectors
\(e_{\rm aux}e_{\rm old}e_\delta\).
They commute with all depth traces, ordinary projectors,
away Hecke operators, the \(q\)-complexes, and the transpose
involution with its transformed fourth-test character.
The actual finite free ordinary module, trace/norm
specialization and finite Hecke-image construction of
AC-G78--AC-G80 therefore apply on this new branch.

\proofstep{4. A compatible new Kummer prime and the actual rank-three seed.}
Keep \(E=\mathbb Q(\sqrt5)\), whose class number is one, and
let \(u\) be the square of a fundamental unit.
Put \(M=E(\mu_p,u^{1/p})\) and
\(\mathcal K=\operatorname{Gal}(M/E(\mu_p))\), of order one
or \(p\).  The image of \(G_{F_{\rm pre}}\) contains
\(\mathcal K\).  Indeed the degree-four argument applies
before \(T_\Theta\) is adjoined, and \(F_1M/F_1\) is unramified
at \(q_{\rm aux}\), whereas \(F_1T_\Theta/F_1\) is totally
ramified of degree \(p\) there.  Thus adjoining \(T_\Theta\)
does not remove this Kummer image.

Choose \(h\in I_5\cap G_{F_{\rm pre}}\) with
\(\eta_E(h)=-1\).  The entire new branch is unramified at
\(5\), so its normalized trace satisfies
\(\operatorname{tr}(ha)=\operatorname{tr}(a)\) for every
group element \(a\), not only at the seed.
Here \(c_p(h)=\chi(h)=\psi_0(h)=1\).
The unchanged coefficient-prime inertia contains \(t\) with
\(c_p(t)=-(1+p)\), \(\chi(t)=1\), and \(\eta_E(t)=1\).
The odd residue \(133\) gives the same negative analytic target
\(-((1+W)+(1+W)^{-1})\) as in the odd branch construction.
Both \(h,t\) act by inversion on \(\mathcal K\).
Conjugating \(h\) by \(n\in G_{F_{\rm pre}}\) whose Kummer
image is \(a\) changes \((ht)^2\) to \(a^4(ht)^2\).
As \(4\) is invertible modulo \(p\), choose \(a\) making
this square one.  Cyclic invariance of the trace preserves
the exact branch scalar identity under this conjugation.

Apply the actual finite-quotient Chebotarev construction with
this corrected element, including \(F_{\rm pre},M\), the
residual characters and \(c_p\bmod p^2\).
It gives a new rational prime \(q\), split in \(F_{\rm pre}\),
inert in \(E\), with
\[
 \chi(q)=\psi_0(q)=1,\qquad v_p(q+1)=1,\qquad
             u\in\mathbb F_{q^2}^{\times p}.
 \tag{G132.17}
\]
Every field used in this compatibility check was fixed before
\(q\).  The actual ramified-disc root argument, or its
fixed-precision normalization version AC-G84, gives the
full-chain point in the finite Hecke image.  The
arbitrary-base-change Cayley--Hamilton support argument puts
it in the actual ordinary weight specialization.
The four tests and the projective
\(\mathcal O[\mathrm{GL}_3(\mathbb F_q)]\)-lattice argument
give an actual finite type
\(\operatorname{Ind}_{P_{2,1}}^{\mathrm{GL}_3(\mathbb F_q)}
 (C_\vartheta^{(2)}\otimes1)\).
Here \(\vartheta\) is regular of order \(269\).
Generic spherical constituents contribute zero to
\(m_{\rm St}-m_1\), nonspherical ones contribute nonnegatively,
and the obtained smooth Steinberg contributes positively.
The unchanged residual genericity witness is available on
this branch by varying the cubic character while fixing the
rational characters.  Thus the Brauer identity produces a
constituent in the retained projector, not after a new average.

Its local parameter at \(q\) is a tame irreducible
two-dimensional block plus an unramified line, with \(N=0\).
The residual Frobenius roots are \(\{-1,1,-1\}\).
The two-dimensional block has roots \(\pm a\), so its
remaining line reduces to \(-1\).
The global residual triple has three distinct characters;
only its middle character is individually conjugate self-dual,
since \(\bar\gamma_*^2\ne1\).
A global \(2+1\) conjugate-self-dual decomposition would
force its line to reduce to that middle \(1\), a contradiction.
Three global lines cannot contain the local irreducible block.
This proves actual global rank-three absolute irreducibility.

The finite-torus character fixes the exact \(p\)-inertia
\((\chi^j,1,\chi^j)\).
The equal endpoint inertia characters have slope gap \(2fw\),
which is not \(f\); the middle inertia differs by the
nontrivial cubic character.  Hence \(N=0\).
The parallel weak-admissibility argument of AC-G79 supplies
the ordered flag with characters
\[
 u_1\chi^j c_p^{133},\qquad u_2,\qquad
 u_3\chi^j c_p^{-133},
 \tag{G132.18}
\]
where the \(u_i\) are unramified.
The actual away type is \((\psi_0,1,\psi_0^{-1})\), with
\(N=0\), by the from-the-start old-type projector.

\proofstep{5. Kill the new type and transport the normalized rank.}
The ray construction of AC-G83 over \(E\), using (G132.17),
gives a totally real dihedral field \(D_q/\mathbb Q\) with
\(D_q/E\) cyclic of degree \(p\), inertia exactly \(p\) at \(q\),
and no other relative finite ramification.
Degree is exactly \(p\) here because \(h_E=1\).
Put
\[
 H_0=F_1T_\Theta D_q,\qquad F_0=H_0^+,\qquad
 H_{\rm ab}\subset H_{\rm ab}T_\Theta\subset H_0 .
 \tag{G132.19}
\]
The two displayed indices are \(p\).
At \(q\), the unramified quadratic extension supplied by
\(E\) splits the tame two-dimensional block, and inertia
index \(p\) kills both characters.  Thus the whole local
parameter becomes unramified.  At \(5\) the old parameter
was unramified before adjoining \(E\).
All other exact old types restrict as stated.
The extension \(H_0/F_0\) is finite-unramified and relatively
split at \(2,3,p\).
Its \(p\)-completions are still unramified degree \(538\):
the unramified local degree from \(D_q/E\) divides \(p\)
and is already contained in the unramified degree-\(2p\)
field from \(F_1T_\Theta\).

Rank-three irreducibility survives the quadratic and
cyclic-\(p\) steps: a nontrivial self-twist of order \(2\)
or \(p\) would have cube one by determinants.
Rank-two irreducibility was already checked on the abelian
field \(H_{\rm ab}T_\Theta\), and persists at its \(p\)-step.
The retained cyclic base change and definite-unitary descent
give the actual rank-three and rank-two automorphic seeds on
\(H_0\), with these flags and types.

Let \(\Omega\) be the original prime-to-\(p\) residual splitting
field, including \(K,\chi,\psi_0,\bar c_p\).
It is unramified at \(5,23\).
The maximal prime-to-\(p\) quotient of the normal host in
(G132.19) is \(H_{\rm ab}\); its three quadratic extensions
over \(K\) ramify at \(5\) or \(23\).
Consequently \(H_0\cap\Omega=K\).
The normal closures of the later selected-place and cyclotomic
extensions have \(p\)-power degree over this host, so this
intersection remains \(K\).
All local and global host degrees have only factors \(2,p\);
in particular no coefficient-prime local field contains
\(\mu_p\), whose degree has the factor \(67\).

Both steps of (G132.19) are unramified at \(p\), hence disjoint
from the totally ramified cyclotomic tower.
Take the full prime-to-\(p\) residual splitting field for
\(\alpha\) at each step.  The commuting finite groups and
the \(p\)-split class modules are precisely those of AC-G100.
Applying its characterwise cyclic rank inequality twice gives
\[
 r_{\alpha,H_0}\le2p^2
        =\tfrac12[F_0:\mathbb Q].
 \tag{G132.20}
\]
This proves the bound on the new actual host.  It invokes no
cubic-sector \(\mu=0\) theorem.

\proofstep{6. Filling, positivity and a common deep host.}
Apply AC-G92 to \(F_0\) and a common prime-to-\(p\) residual
cover containing \(H_0,\mu_p\), the adjacent characters and
the rational sectors
\(\bar\gamma_*^{-1}\eta_K,\bar\gamma_*^{-2}\eta_K\).
The selected-place extension \(T'/F_0\) is cyclic of
\(p\)-power degree and splits every \(p\)-place.
Its normal closure over the old host has \(p\)-power degree.
AC-G93's actual ray-class quotient gives nonzero strict
adjacent and central anti classes, and the two rational
strict sectors.  The latter are odd and nontrivial.
The local invariant groups vanish: the cubic order three
cannot cancel an inertia character of order dividing \(268\),
and the rational endpoint square has order \(134\).
These properties persist after \(p\)-power restriction.

The full dual for the mixed obstruction is the even character
\[
 \lambda=\eta_K\bar\psi_0^2\bar c_p^{2e_*+1}
         =\eta_K\bar c_p^{267}.
 \tag{G132.21}
\]
Its \(p\)-inertia is nontrivial and its splitting field is
totally real and abelian.  After adjoining that field and
\(\mu_p\), the real host is obtained by finite \(p\)-extensions
from an abelian field, including the selected-place extension.
The retained abelian \(\mu=0\) theorem and finite-\(p\) ascent
bound its full \(H^1\) along the cyclotomic tower.
The exact away-kernel/full-local cone comparison of
AC-G90--AC-G93 applies on this possibly non-\(\mathbb Q\)-Galois
selected host.  Restriction in the dual \(H^1\) is injective
and eventually an isomorphism; corestriction is then zero
because its composite with restriction is multiplication by
\(p\).  Duality makes restriction on the degree-two mixed
obstruction zero.  Nonzero adjacent and central classes
persist.  AC-G90 supplies the actual polarized long entry
and polynomial Low paths, not only a vanishing obstruction
class.

The selected extension splits \(p\), so it is disjoint from
the cyclotomic tower.  A further application of AC-G100
preserves the ratio \(r_\alpha/d\le1/2\).
The rational central and trivial sectors have bounded
dimensions by finite-\(p\) ascent from their abelian splitting
fields.  The adjacent inertia is nontrivial, so the local
permutation correction in AC-G74 is zero.
AC-G98 gives, on successive cyclotomic layers,
\[
 a_n=r_\alpha p^n+O(1),\quad h_n=O(1),\quad z_n=O(1),
 \quad d_n=dp^n,\quad
                 h_n+a_n+z_n+1<d_n
 \tag{G132.22}
\]
for all sufficiently large \(n\).
Choose one layer beyond this threshold, the mixed-filling
threshold and all local \(d_v>5\) thresholds.
On it the rational frame bounds are also satisfied, as
checked next.  Denote the resulting CM host by \(H/F\).
The seeds remain irreducible through the cyclic-\(p\) steps
and give actual cuspidal base changes with the same finite
types.  No effective depth or discriminant bound is asserted.

\proofstep{7. Reconstruct the rational corner and its native support.}
We give the changed-residual checks needed to replay
AC-G106-native and AC-G121, rather than assume their
residual-exponent-\(267\) instances.
The rational splitting fields for \(\bar\gamma_*\) are abelian
over the initial \(H_{\rm ab}T_\Theta\).
The abelian Iwasawa theorem and finite-\(p\) ascent make their
strict and trivial dimensions bounded on the tower.
The positive rational sectors of Step 6, the exact Euler and
\(\Phi\)-finiteness statement AC-G96, and the full-cyclotomic
kernel inertia generators and finite class exponent AC-G99
give every corresponding rational frame field.
On inertia \(\bar\gamma_*^2\) has order \(134\), and
\(1\pm2e_*\equiv267,3\pmod{268}\) are nonzero.
Thus the rational \(\mathrm{Frame.Ok}\) and
\(\mathrm{Frame.Big}\) inequalities hold on a sufficiently
deep common layer already permitted in Step 6.

The scalar auxiliary place is actual.
The conductor-\(4\) field is disjoint from the joint cubic
and conductor-\(p\) fields by ramification at \(2\).
Choose their compatible element with
\(\chi=1,\psi_0=-1,\bar c_p=-1\); it fixes \(K\) and has
\(\bar\gamma_*=1\).
The intersection calculation in Step 5 extends it fixing
the host normal closure.  Chebotarev supplies a fresh
completely split rational prime and a residue-degree-one
place \(v_0\) with these values.
The native residual characteristic polynomial is scalar
there.  The residual Sylvester operator for \(\mathrm{N4}\)
has only the eigenvalue
\(\mu(1-\bar c_p(v_0))\), which is nonzero; it is therefore
invertible even when the residual matrix has a nilpotent part.
This is the determinant-only \(\mathrm{N4}\) proof, with
no scalar-matrix assumption in characteristic zero.

Take the literal \(\mathrm{HidaSetting}\), with \(S\) containing
the places over \(2,3\), and retained \(\mathrm{LevEx}\).
Its \(\mathrm{p\_unr}\) concerns the new auxiliary rational
prime, not \(269\).  All coefficient-prime local maps are
the full decomposition groups.
Since their fields have no \(\mu_p\), local reciprocity
identifies the source's \(p\)-primary torsion preimage
\(\mathrm{mu}\) with the closed commutator; \(\mathrm{th}=1\)
is valid.  Enlarge the host coefficient field to contain its
embeddings, \(\kappa_2\), and the seeds, as allowed by
\(\mathrm{Coef}/\mathrm{Large}\).  This makes no change to
\(\mathcal O_0\).

An actual split residual rank-two lattice is available.
Diagonalize \(V(c)\) to \(\operatorname{diag}(-1,1)\).
The distinct residual characters give an invariant residual
line, so one off-diagonal entry family is divisible by a
uniformizer.  After adjoining its square root, scale one
basis vector so both off-diagonal families are divisible
by the new uniformizer.  The diagonal \(c\)-matrix is
unchanged and the residual representation is split.
On \(G_H\) the determinant is one and
\[
 V((cgc^{-1})^{-1})=J_2V(g)^tJ_2,\qquad
             J_2=\begin{pmatrix}0&1\\1&0\end{pmatrix}.
 \tag{G132.23}
\]
The difference is diagonal, giving the required normalization.
The exact away type is \(\psi_0\oplus\psi_0^{-1}\), now that
\(\Theta\) is killed.  It kills the normal closure \(N_\gamma\)
of all \(I_w\cap\ker\bar\gamma_*\), \(w\nmid p\).
Thus it is an actual normalized integral representation of
\(\Gamma_\gamma=G_F/N_\gamma\), with the full ordered local
flags.  This is the host substitute for the old rational
\(\mathrm{Dat.Kill}\); no original point is used.

The literal rational \(\mathrm{space\_ex}\) constructs
\(B_2\subset A_2\subset R_2\) representing the full complete
local normalized ordinary functor.
Use the actual conductor-three \(\kappa_2\) of AC-G106-twist,
with \(\kappa_2\kappa_2^c=c_p^{-1}\), on the native group.
The native CM kernel maps onto
\(G_\gamma=\ker(\eta_K:\Gamma_\gamma\to\{\pm1\})\),
preserving conjugation, the full cyclotomic character and
all the full local maps.  The character \(\kappa_2\) need
not descend further to \(G_\gamma\).
The representation \(\kappa_2r_2\), local characters
\(\kappa_2\lambda_i\), and residual characters
\(\bar\kappa_2\bar\gamma_*^{\,1-2i}\) define the native Pair.
For its inertia field, the joint residual kernel can include
\(\bar\gamma_*=-1\): a finite prime-to-\(p\) order universal
matrix with scalar residue equals its Teichmuller scalar,
by factoring \(X^m-1\) and inverting the other factor even
over nonreduced rings.  The finite away \(\kappa_2\) value
cancels that scalar.  Hence the literal \(\mathrm{Pair.I}\)
holds.  The determinant law is coefficientwise in \(B_2\)
and continuous by the weight-zero atom bound.

The native map \(\Phi_2:\mathcal B_2\to B_2\) is surjective:
its closed complete image contains the traces and local
characters.  For the normalized
\(Q=r_2(g_0)-r_2(g_0^\star)=\operatorname{diag}(a,-a)\),
\(\operatorname{tr}(Q^2)/2=a^2\) is in that image and
\(\bar a\ne0\).  Hensel lifting the specified root puts \(a\)
there.  Coordinate projectors and their word traces then
give all diagonal entries and all
\(r_{01}(g)r_{10}(h)\), which generate the completed
weight-zero ring.

The actual arithmetic point uses the unitary base change
\(\pi_0\) of \(V\), and the fixed Hecke character \(Y\)
corresponding to \(\kappa_2\):
\(\Pi=\pi_0\otimes(Y|\cdot|^{1/2})\circ\det\).
The identity \(YY^c=|\cdot|^{-1}\) proves conjugate
self-duality; the geometric \(q^{1/2}\) attachment factor
leaves exactly \(V\kappa_2\).
Its regular algebraic infinity pair is
\((133-\tfrac12,-133-\tfrac12)\), and the ordered local
weights are the two weights of \(\kappa_2\) shifted by
\(-133,+133\).  Their gap and the finite phase verify the
literal arithmetic \(\mathrm{PT}\) clauses.
Thus retained \(\mathrm{ModU}\) supplies a proper new \(J_2\).
Native finite/torsionfree control and the actual weight map
give its component dimension \(2d+1\).
The exact native \(\mathrm{Setup}\) and full actual seed point
give \(\mathrm{SeedIn}\): the whole native component maps
back to \(B_2\), with composite equal to its native quotient,
by the \(\mathrm{seedOf}\) construction used in AC-G106-native.
Thus this assertion is stronger than having one closed point.
Together with actual finite-weight specialization and the
rank-two component theorem it gives the uniform
off-triangular ideal of AC-G106-native.
Transport by AC-G106-even, retaining the independent middle
character, gives the actual \(M_E\) for the new tuple.
It is not an assumed ideal from the old e267 tuple.

\proofstep{8. Reconstruct cubic support and propagate it.}
Form the cubic quotient by the joint residual away inertia
kernel for (G132.3).  All local separating characters and
the normalized inertia witness have the required distinct
values: cube the local witness to kill \(\chi\) while
retaining the order-\(268\) rational separation.
AC-G94 constructs the full universal ring \(R\), with
labelled Hensel projectors and ordered-product relations
over arbitrary complete local test rings.
AC-G95 constructs the actual graded \(A\) and complete
invariant \(B\); AC-G96 gives the arbitrary-module Euler
and finiteness conditions, and AC-G99 gives the full
local generator and class-exponent conditions.
The filled Low points and paths in Step 6 are actual
points of this chart.

For rank three use \(\kappa_3=c_p^{-1}\).
The universal \(\kappa_3r\), all three local characters,
the determinant law in \(B\), and the joint-kernel
inertia check give the literal new native Pair and Resid.
The finite away type is prime to \(p\), so reduction
identifies its Teichmuller constants; at \(p\) the
literal \(\mathrm{mu}/\mathrm{th}\) clauses are those
verified in Step 7.
The trace-projector argument AC-G102 makes
\(\Phi_3:\mathcal B_3\to B\) surjective.
The actual rank-three seed of Steps 4--6 gives the
arithmetic native point with ordered weights
\[
                  (1-133,1,1+133).
 \tag{G132.24}
\]
Its finite \((\chi^j,1,\chi^j)\) phases are retained.
The weights are strictly increasing and no character
ratio is \(c_p\).  The actual \(\mathrm{PT}\) implication
makes \(J_3\) proper.  Native finite/torsionfree weight
control, nonzero support, and the actual \(\mathrm{LamEx}\)
map give dimension \(3d+1\), with the unramified and
\(\mathrm{N4}\) fields as above.

AC-G107 applies to this seed and these universal data:
its entire component, not just its closed point, maps
back by \(\beta:B\to C\) with \(\beta\Phi_3\) equal to
the native quotient.  It gives the strong seed prime
of dimension \(3d+1\) containing \(I\).
There is no assertion that \(B/I\) is torsionfree.
AC-G102 supplies actual finite weight support, AC-G104
the residually normalized native prime conditions,
AC-G101 the local presentations and all-open
character-separation loci, and AC-G105 the parity
presentation and chaining.
The filled lower paths of AC-G90/G95, the strict
dimension inequality (G132.22), the strong seed, and
the actual uniform \(M_E\) of Step 7 are now all present.
The full composition theorem AC-G108 gives exactly
(G132.4)--(G132.5).
The first ideal \(I\) is used for Type I points;
the intersected \(I_{\rm off}\) is used for curves.

All these host rings, residual records and support
ideals depend on the fixed \(\bar\gamma_*\).
For the core progression \(e\equiv133\pmod{268}\)
this residual character is constant.  The one seed at
\(e_*=133\) therefore suffices for every member.
The full ordinary functor leaves its integral local
characters variable; it does not force every point
to have the seed's characteristic-zero weight.

\proofstep{9. Original curves, rigidity and integral lengths.}
Take an actual original characteristic-\(p\) DVR curve.
Use the AC-G109 residue extension
\(k'[[z]]\hookrightarrow k''[[z]]=V'\), where \(k''\)
contains the host residue field \(k\), and the host
test ring \(C=k+zV'\).
The original \(\mathcal O_0\) and \(S_j\) are unchanged.
Exact away-kernel killing and the full flags of the
original selected component AC-G65-flag give the
restricted representation over \(C\).
The paired Hensel projectors and strict orthogonal
normalization give an actual map \(R\to C\).
Its inertia characters reduce to (G132.3), so it kills
the actual \(\lambda_\theta\).

If the curve were off both triangular loci, (G132.5)
would kill its positive and negative balanced monomials
in \(\mathfrak m_B\), contradicting any surviving
opposite-sign entries in the domain \(C\).
It is triangular.  Removing the residual diagonal
characters leaves everywhere-unramified principal-unit
characters with finite class-group exponent, hence one
in characteristic \(p\).
Some power of each original group element lies in
\(G_H\); its eigenvalues therefore satisfy finite-field
polynomials.  Their prescribed reductions identify
the original constant characteristic coefficients,
also when the host index is divisible by \(p\).
The generic shape/trace-pairing argument and integral
finite criterion of AC-G109 give original M2.

For an original Type I characteristic-zero point take
a finite value field containing it and the host
coefficients, and the inverse image of \(k\) in its
integers as host test algebra.
The same exact flags and away conditions give its
actual point of \(A\).
AC-G110's all-open and cyclotomic-kernel irreducibility
arguments apply to this separated tuple.
The first line of (G132.5) makes its native point kill
\(J_3\).  Native \(\mathrm{end2}\) gives automorphy
with that point's own weights \((1-e,1,1+e)\).
The retained polarized adjoint theorem and
restriction/corestriction give original adjoint
rigidity.  The cohomology here has characteristic-zero
coefficients, so a \(p\)-divisible host degree is
invertible.

AC-G117 excludes every original Type II point under
the retained original coefficient and modularity
hypotheses.  This does not exclude the auxiliary
seed used above, which has a different rational level
and is not an original \(\mathrm{Dat.IsPt}\).
Every two-sided original point is consequently Type I
and rigid.  The exact normalized dual-point argument
of AC-G122 puts both cobounding derivations in the
one line spanned by \(\operatorname{diag}(1,0,-1)\);
it gives literal M4.  AC-G66 gives M1 and M3.
At the degree-two place over \(11\), both the
quadratic and cubic characters are one, so the seed
trace is \(b>3\).  Actual Type I traces have complex
absolute value at most three.  The degree-zero trace
difference gives M5.

The host is reached from \(K\) by cyclic steps of
degree \(2\) or \(269\).  AC-G111 gives its exact
original-field descent and witness.
The original inertia, not the auxiliary seed level,
gives principal compact level \(269N\) by AC-G112.
Thus the classifier contains neither \(q_{\rm aux}\)
nor the later ray prime in its conductor bound.
The original \(\mathrm{LemmaM},\mathrm{LemmaM'}\) and
strict extraction apply with
\(v_{\mathcal O_0}=e_{\mathcal O_0}v_{269}\).
They give (G132.8); the trace-norm construction on
the common original-field classifier gives (G132.9).
In particular
\(\#J_0C_1Y^{A_1}(2e\log121+1)
 \le5C_1\#J_0Y^{A_1+1}\log11\).
No host-coefficient ramification index rescales
these original lengths.

\proofstep{10. The complementary four original families.}
Each discriminant in (G132.6) is fundamental:
it is \(4d\) with \(d\) squarefree and \(d\equiv3\pmod4\).
The four characters are odd quadratic, their conductors
are the stated \(N_{s,t}\), and
\(\operatorname{comp}_p\psi_{s,t}=\tau^t\), \(\tau=\chi_{269}\).
The exponents are positive odd and
\(2e\equiv266\pmod{268}\).
The parity, genericity, prime-to-\(p\) character-order
and even-modulus clauses of the original GD hold.
Its prime-to-\(p\) character/cyclotomic field is
unramified outside \(2,5,p\), so cannot contain
the field \(K\) ramified at \(3\); \(3\nmid N\) and
\(p\nmid3h_K\).  The actual marked and selected-local
data remain the stated original inputs for each row.

The field \(E=\mathbb Q(\sqrt5)\) is already in \(F\),
so \(\chi_5|_{G_F}=1\), while
\(\bar\tau=\bar c_p^{134}\).  Hence on the full host
group and all its full local groups
\[
 \bar\psi_{s,t}\bar c_p^{e}
       =\bar\psi_0\bar c_p^{e+134t}
       =\bar\gamma_* .
 \tag{G132.25}
\]
This identifies the actual host residual record and
the full ordinary functor already constructed above.
It does not assert a new disjointness statement for
the enlarged original splitting field, which may
contain \(E\).
On away inertia \(\tau\) is unramified and \(\chi_5\)
is trivial on the host.  In particular at \(5\) the
original modified Kill condition kills all host
inertia.  At every other away place it gives exactly
the same joint kernel.  Thus the maps to the host
chart, not a map between original deformation rings,
have the required away domain.

At \(p\), \(126^2\equiv5\pmod{269}\) with nonzero
derivative, so \(\chi_5|_{G_{\mathbb Q_p}}=1\).
For \(t=0\) the two rows have identical full local
data at a fixed weight.  For \(t=1\) apply its own
AC-G65-local and AC-G65-flag theorem to
\(P_{*,s,1,e}\), retaining the quadratic phase.
It gives the strict full \(G_{K_p}\)-stable flag
over every test algebra, with ordered inertia
\[
 \chi^j\psi_{s,t}c_p^e,\quad1,\quad
                      \chi^j\psi_{s,t}^{-1}c_p^{-e},
 \tag{G132.26}
\]
finite type \((\chi^j\tau,1,\chi^j\tau)\), and
weights \((-e,0,e)\).
The hypotheses \(\psi_{s,t}^2|_{I_p}=1\), \(e\ge1\),
and \(268\nmid2e\) hold.  No arbitrary ordinary
component is substituted.
The reductions of (G132.26) are exactly (G132.3).
On the full finite unit torus, the new algebraic
weight and quadratic phase change the endpoint
factors by exponents
\[
                  134t+(e-133),\qquad134t-(e-133),
 \tag{G132.27}
\]
both divisible by \(268\).
Thus the exact finite character is (G132.16).
The integral local characters and characteristic-zero
types remain those of their own original models.

In characteristic \(p\), (G132.26) kills
\(\lambda_\theta\) after host restriction, giving
the same original M2 argument of Step 9.
In characteristic zero, its finite phase supplies
the native \(\mathrm{IsWtPart.ph}\) condition and
its weights become \((1-e,1,1+e)\).
Here \(e\ge133\); no ratio can equal \(c_p\), since
its weight difference has absolute value \(e\)
or \(2e\).  The same actual \(I\), not just
\(I_{\rm off}\), gives native Type I automorphy
and rigidity.  AC-G117 explicitly permits the
ramified quadratic \(p\)-component and excludes
original Type II in all four rows.
Thus M1--M5 hold for every row.
The same Frobenius at \(11\) is prime to every
\(3pN_{s,t}\) and has the same seed trace \(b\).
AC-G112-local permits \(p\mid N\) and gives
the uniform principal level \(pN_{s,t}\).
Repeating the original valuation argument with
\(Y=N_{s,t}(e+1)\) proves (G132.7)--(G132.9)
for all four rows and completes the proof.
\end{proof}

\paragraph{Scope, geometric dictionary, and source accounting.}
The class input and the entire host/native composition for e133
are new here.  In particular this theorem never identifies its
\(J_2,J_3\) with the e267 ideals of AC-G123.
The four-row extension uses that new support and individually
checks the original selected local components; it has no separate
endpoint-existence or support premise.
It gives no quantitative first cyclotomic layer, all-character
or all-coefficient-prime assertion, or the original H04.4 counting
budget.  The strict tower certificates at a fixed prime alone do
not put their characters in an unrelated prescribed growing host.

The auxiliary seed lies in the actual tame modular-curve cusp
cohomology of AC-G119.  Restriction to the displayed finite CM
fields and tensoring by the specified characters have the already
audited finite-cover/local-system interpretation of SGA~1,
Expose~V, sections~4 and~7.  The host point is a point of the
normalized ordinary formal deformation chart; comparison with
an original model is point by point or curve by curve.
No new motive or cohomology-comparison theorem is claimed.

The applied complete AC-G120, AC-G121, AC-G123 and AC-G124 proofs
were reread for this changed tuple.  The exact AC-G103
Kummer/Shapiro and class-sector statements, AC-G115, AC-G126,
the successful D6187 certificate and its flag-one specification,
and the complete AC-G131 proof and independently replayed
evidence are retained with exact hashes in the source matrix.
The unchanged primary theorem scopes and complete native-file
audits recorded in AC-G94--AC-G112 and AC-G119--AC-G124 remain
explicit inputs.  This composition does not claim a new proof
of those deep inputs or a Lean build.
